\documentclass[11pt]{article}
\usepackage[T1]{fontenc}
\usepackage{lmodern,amsmath,amssymb,amsthm,graphicx,booktabs,tabularx,microtype}
\usepackage[letterpaper,margin=0.8in]{geometry}
\usepackage[hidelinks]{hyperref}
\newtheorem{proposition}{Proposition}
\setlength{\parindent}{0pt}
\setlength{\parskip}{5pt}
\setlength{\emergencystretch}{2em}
\allowdisplaybreaks
\begin{document}
% Append immediately after comparison_bounds_appendix.tex, within Appendix A.
% Uses the existing construction labels and bibliography. Requires graphicx.
\providecommand{\qeqref}[1]{(\ref{#1})}


\section{Selection of the relative bandwidth}
\label{sec:qi-bandwidth}

We derive the tanh bandwidth prescription from the pole contribution in
the representation proof's quadrature decomposition. We retain the geometry
$h=2/N$, $\gamma=\lambda/h$, $d=\pi/(2\gamma)$, and
$W=N+2R+1$, with $0<\lambda\le1$. In particular,
\begin{equation}
 \frac{2\pi d}{h}=\frac{\pi^2}{\lambda},
 \label{eq:bandwidth-exponent}
\end{equation}
so the exponential factor supplied by the midpoint bound is
$e^{-\pi^2/\lambda}$, as in the representation theorem's pole estimate.

% Insertion for the supplied bandwidth appendix.
% Place after eq:bandwidth-exponent and before "Frequency-dependent pole bound".
% Retain the supplied tanh proposition and proof, compact-score comparison,
% frequency-estimation discussion, and tolerance qualifications.
% Requires amsmath, amssymb, booktabs, graphicx, tabularx.
\providecommand{\BandwidthWidthFigure}{figures/bandwidth_width_rules.pdf}
\providecommand{\BandwidthErrorFigure}{figures/bandwidth_error_rules.pdf}

\subsection{The general rule and its tanh default}
\label{app:bw-general-rule}

We first state the bandwidth prescription for an activation $\sigma$ with a
localized derivative $K=\sigma^{(r)}$. We assume $K\in L^1(\mathbb R)$ and
$\widehat K(0)\ne0$, and use the angular-frequency convention
\[
 \widehat K(\xi)=\int_{\mathbb R}K(t)e^{-i\xi t}\,dt.
\]
The general rule selects a bandwidth from the normalized Fourier magnitude
\begin{equation}
 G_K(\lambda):=
 \frac{|\widehat K(2\pi/\lambda)|}{|\widehat K(0)|},
 \qquad
 G_K(\lambda_{\rm general})=\varepsilon_{\rm eff}.
 \label{eq:bw-general-score}
\end{equation}
It depends on the activation and the chosen tolerance, and not on width or
target frequency. Multiplying $K$ by a nonzero constant does not change the
rule. The relevant transforms for the four activations in our figures are
\begin{equation}
\begin{array}{c|c|c}
 \sigma(x)&r&\widehat K(\xi)\\\hline
 \tanh x&1&\displaystyle\frac{\pi\xi}{\sinh(\pi\xi/2)}\\[5pt]
 x\Phi(x)\quad\text{(exact GELU)}&2&
       (1+\xi^2)e^{-\xi^2/2}\\[3pt]
 \displaystyle\frac{x}{1+e^{-x}}\quad\text{(SiLU)}&2&
       \displaystyle\frac{(\pi\xi)^2\cosh(\pi\xi)}{\sinh^2(\pi\xi)}\\[5pt]
 \operatorname{erf}(x)&1&2e^{-\xi^2/4}
\end{array}
\label{eq:bw-three-transforms}
\end{equation}
Here $\Phi$ is the standard normal cumulative distribution function. The
values at zero are defined by continuity: $2$, $1$, $1$, and $2$, respectively.
At $\varepsilon_{\rm eff}=2^{-52}$, the roots on the decaying Fourier tail are
\[
 \lambda_{\rm general}^{\tanh}=0.2440764,\qquad
 \lambda_{\rm general}^{\rm GELU}=0.6985705,\qquad
 \lambda_{\rm general}^{\rm SiLU}=0.4453827.
\]
For erf, $\lambda_{\rm general}=\pi/\sqrt{\log(1/\varepsilon_{\rm eff})}$.
For completeness, GELU has $K(x)=(2-x^2)\phi(x)$, where $\phi=\Phi'$;
the Gaussian transform and differentiation give its table entry. For SiLU,
write $\ell(x)=(1+e^{-x})^{-1}$ and $L=\ell'$. Then
$K=2L+xL'$ and $\widehat L(\xi)=\pi\xi/\sinh(\pi\xi)$, so
$\widehat K=\widehat L-\xi\widehat L'$, giving the stated expression.
For erf, $K=2e^{-x^2}/\sqrt\pi$ gives the last entry directly.
Thus $\lambda=1/4$ is a convenient rounded FP64 tanh default, rather than a
third activation-independent rule. It is not exactly the root of
\eqref{eq:bw-general-score}: the tanh score at $1/4$ is
$5.6511\times10^{-16}$.

The quarter-scale default also has an independent justification under the
admissibility hypotheses of the tanh representation theorem. If $B$ bounds
the analytic target on the prescribed complex neighborhood, that proof gives
\begin{equation}
 E_{\rm alias}\le4B\mathcal U(\pi^2/\lambda),\qquad
 \mathcal U(a)=\frac{e^{-a}}{(1-e^{-a})(1-e^{-2a})}.
 \label{eq:bw-quarter-class-bound}
\end{equation}
This uses the representation theorem's condition $\gamma\delta\ge4\pi$,
where $\delta$ is the analytic-neighborhood radius.
At $\lambda=1/4$, $4\mathcal U(4\pi^2)=2.8629\times10^{-17}$.
This statement concerns the aliasing contribution; the boundary, resolution,
and numerical-recovery contributions still require their own estimates.

\subsection{Why the kernel ratio appears}
\label{app:bw-replicas}

Use the grid coordinate $u=(x+1)/h$, where $h=2/N$. A finite lattice sum of
derivative kernels has Fourier transform
\begin{equation}
 \widehat{\sum_j a_jK(\lambda(u-j))}(\theta)
 =\lambda^{-1}\widehat K(\theta/\lambda)C(\theta),\qquad
 C(\theta)=\sum_j a_je^{-ij\theta},\quad
 C(\theta+2\pi m)=C(\theta).
 \label{eq:bw-lattice-spectrum}
\end{equation}
The coefficient factor repeats at every lattice frequency. Dividing a
replica by the retained Fourier component therefore cancels that factor.
For a nonzero mode $\theta=h\omega$, integrating $r$ times gives the output
replica ratio
\begin{equation}
 t_m(\lambda,\theta)=
 \left(\frac{\theta}{\theta+2\pi m}\right)^r
 \frac{\widehat K((\theta+2\pi m)/\lambda)}
      {\widehat K(\theta/\lambda)},\qquad m\ne0.
 \label{eq:bw-output-replica}
\end{equation}
The central transform must be nonzero. Polynomial integration constants of
degree below $r$ are separate from these nonzero-frequency ratios.
Equation~\eqref{eq:bw-output-replica} explains both ingredients of the
refinement: the shifted Fourier arguments and the frequency factors from
integration. It does not require the fitted readout coefficients.

\subsection{The frequency-dependent refinement}
\label{app:bw-refined-rule}

For $0<\theta<\pi$, retain the nearest replica pair and define
\begin{equation}
 T_{K,r}(\lambda,\theta)=
 \frac{
 \left(\frac{\theta}{2\pi-\theta}\right)^r
   |\widehat K((2\pi-\theta)/\lambda)|
 +\left(\frac{\theta}{2\pi+\theta}\right)^r
   |\widehat K((2\pi+\theta)/\lambda)|}
 {|\widehat K(\theta/\lambda)|}.
 \label{eq:bw-first-pair}
\end{equation}
Given a representative angular frequency $\bar\omega$, the refined rule is
\begin{equation}
 T_{K,r}(\lambda_{\rm refined},2\bar\omega/N)
 =\varepsilon_{\rm eff}.
 \label{eq:bw-refined-root}
\end{equation}
Both prescriptions select the largest admissible bandwidth within a stated
search interval. A bracketed root solve is used on an increasing branch;
if the entire interval meets the budget, its upper endpoint is returned,
and an infeasible interval is reported rather than extrapolated.

The general score retains the rapidly varying Fourier decay while dropping
the width- and target-dependent prefactors. More precisely, at fixed
$\lambda$ and $\theta\downarrow0$,
\begin{equation}
 \frac{T_{K,r}(\lambda,\theta)}{2(\theta/(2\pi))^r}
 \longrightarrow G_K(\lambda).
 \label{eq:bw-general-refined-limit}
\end{equation}
The division by $\theta^r$ matters: the refined score itself does not tend
to the general score. Nor does this limit imply that roots at a fixed
tolerance converge to $\lambda_{\rm general}$ as width tends to infinity.

The function enters this scalar rule only through $\bar\omega$.
Consequently its width curves are exact horizontal rescalings in $N$:
if $\Lambda_{K,r}(z)$ solves $T_{K,r}(\lambda,2/z)=\varepsilon_{\rm eff}$,
then $\lambda_{\rm refined}(N,\bar\omega)=\Lambda_{K,r}(N/\bar\omega)$.
Converting to total width $W=N+2\lceil\sqrt N\rceil+1$ slightly changes this
rescaling. It introduces no additional target dependence.

The reason a constant general rule can remain useful is threshold
sensitivity, rather than an asymptote of the refined roots. In the
small-angle regime, put $C_\theta=2(\theta/(2\pi))^r$. If $\lambda_b$ solves
$G_K(\lambda_b)=\varepsilon_{\rm eff}$ and $\lambda_s$ solves the surrogate
$C_\theta G_K(\lambda_s)=\varepsilon_{\rm eff}$, the mean value theorem gives
\begin{equation}
 \log\frac{\lambda_s}{\lambda_b}
 =-\frac{\log C_\theta}{D_K(\widetilde\lambda)},\qquad
 D_K(\lambda)=\frac{d\log G_K(\lambda)}{d\log\lambda},
 \label{eq:bw-threshold-sensitivity}
\end{equation}
for some $\widetilde\lambda$ between the two roots. For tanh,
$D_K=a\coth a-1$, $a=\pi^2/\lambda$; its value at the FP64 general root is
about $39.44$. A large change in an algebraic prefactor can therefore move
the selected bandwidth much less. This comparison requires both a small
central angle and negligible shifts in the Fourier tail: for tanh,
$\pi\theta/\lambda\ll1$; for GELU, $2\pi\theta/\lambda^2\ll1$ and
$\theta/\lambda\ll1$. The full refined score retains these shifts instead.

For tanh, set $a=\pi^2/\lambda$ and
$s=\pi\theta/(2\lambda)$. Direct substitution gives
\begin{equation}
 T_{K,1}(\lambda,\theta)
 =\sinh s\,[\operatorname{csch}(a-s)+\operatorname{csch}(a+s)]
 \simeq2e^{-a}\sinh(2s).
 \label{eq:bw-tanh-reduction}
\end{equation}
With $\theta=2\bar\omega/N$, the final expression is exactly the compact
score used in \eqref{eq:bandwidth-prescription} below. The supplied
comparison \eqref{eq:bandwidth-score-comparison} controls this reduction.
Thus the first-pair rule for a general activation and the compact tanh rule
are two forms of the same refinement, with the latter using a controlled
tail approximation.

For an even positive transform whose logarithm is concave on the alias
tail, put
\[
 \rho_K=
 \frac{\widehat K((4\pi-\theta)/\lambda)}
      {\widehat K((2\pi-\theta)/\lambda)}<1.
\]
The complete unanchored replica sum satisfies
\begin{equation}
 T_{K,r}\le\sum_{m\ne0}|t_m|
 \le\frac{T_{K,r}}{1-\rho_K}.
 \label{eq:bw-tail-control}
\end{equation}
Successive transform ratios are at most $\rho_K$ by log-concavity, and the
integration factors also decrease, so summing the two geometric tails
proves the upper inequality. The finite-contour tanh construction has an
additional factor two from anchoring, as proved below. The general replica
identity alone does not establish that finite-domain theorem for GELU or
SiLU or erf. A representative frequency is also a practical reduction: for a
multi-frequency target, a certified alias bound retains the full weighted
spectrum, as in \eqref{eq:bandwidth-weighted-bound}.

\subsection{Inputs and evaluation of the figures}
Both figures use $\varepsilon_{\rm eff}=2^{-52}$, the spacing above one in
binary64. The scalar solve uses the exact first-pair score, without the
anchoring factor or later-pair correction; these are practical selections,
not total-error certificates. Every displayed root is bracketed on
$[0.003,1.5]$. The width figure displays $32\le N\le8192$ on a linear
$N$ axis. A curve is omitted wherever its representative frequency
violates $2\bar\omega/N<\pi$. Its frequency inputs are:
\begin{center}
\small
\begin{tabularx}{\linewidth}{l r X}
\toprule
Target & $\bar\omega$ & Frequency convention\\
\midrule
$\sin(4\pi x)$ & $4\pi$ & Exact single frequency.\\
$\sin(24\pi x)$ & $24\pi$ & Exact single frequency.\\
$\sin(2\pi x)+10^{-3}\sin(40\pi x)$
 & $2.04\pi/1.001$ & Exact amplitude-weighted line centroid.\\
$\sin(8\pi(x+1)^2)$ & $16\pi$ & Mean absolute phase derivative on $[-1,1]$.\\
$1/(1+25x^2)$ & $5$ & Whole-line Fourier-amplitude centroid.\\
$1/(1+100x^2)$ & $10$ & Whole-line Fourier-amplitude centroid.\\
$\log(\cosh(10x))/10$ & $6.45255$ & Fixed 4096-point period-two DFT amplitude centroid.\\
$e^{-8(x-0.2)^2}\sin(8\pi x)$ & $8\pi$ & Carrier-frequency approximation.\\
\bottomrule
\end{tabularx}
\end{center}
The log-cosh transform uses $x_\ell=-1+2\ell/4096$, omits the repeated
endpoint, retains DC in the centroid denominator, and applies no window or
filter. Although endpoint values agree, the periodic extension has a
derivative jump; this estimate depends on transform resolution, which is
held fixed as width changes. The chirp and packet inputs are characteristic
frequency estimates, not exact amplitude centroids.

The error figure uses four different targets:
$\sin(2\pi x)+\tfrac12\sin(6\pi x)+\tfrac14\sin(10\pi x)$,
$e^{-20x^2}$, $e^{\sin(3\pi x)}$, and $\operatorname{sech}(5x)$.
Their frequency inputs are, respectively,
$30\pi/7$, $\sqrt{80/\pi}$, $6.349190385$, and
$40\mathsf C/\pi^2$, where $\mathsf C=\sum_{j\ge0}(-1)^j/(2j+1)^2$
is Catalan's constant. The first and third use discrete Fourier-amplitude
centroids (including DC); the second and fourth use whole-line transform
centroids. In particular,
$\widehat{\operatorname{sech}(5\cdot)}(\omega)
=(\pi/5)\operatorname{sech}(\pi\omega/10)$ gives the final value.

All fits use uniformly spaced centers
$c_j=-1+2j/N$, $-32\le j\le N+32$, and
$N\in\{32,48,64,96,128,256\}$.
Thus $W=N+65$ and the halo is fixed at 32 per side, matching
the archived experiment. Each feature matrix includes a constant column;
there is no separate affine skip. Readouts are fit on 2003 uniform points
by FP64 SVD least squares with relative singular-value cutoff $2^{-52}$.
Relative $L^2$ errors are evaluated on 4001 uniform points. The displayed
curves combine 1440 archived observations and 2400 added observations under
this protocol; 96 additional fits evaluate the exact refined selections.
No curve smoothing is applied. The network fits test the practical rule;
they do not instantiate the boundary-corrected tanh construction in the
supremum-norm theorem.


\subsection{Frequency-dependent pole bound}
For a positive angular frequency $\omega$, put
\begin{equation}
 \begin{gathered}
 A_\lambda=\frac{\pi^2}{\lambda},\qquad
 t_\omega=\omega d,\qquad
 \chi_\omega=2t_\omega=\frac{2\pi\omega}{N\lambda},\\
 \mathcal T(\lambda,h\omega)=\sinh t_\omega
 \left[\frac1{\sinh(A_\lambda-t_\omega)}
       +\frac1{\sinh(A_\lambda+t_\omega)}\right],
 \end{gathered}
 \label{eq:bandwidth-pair}
\end{equation}
with $\mathcal T(\lambda,0)=0$. The first-pair score bounds the complete
pole contribution as follows.

\begin{proposition}[Frequency-dependent aliasing bound]
\label{prop:bandwidth-alias}
For a finite exponential sum $f(x)=\sum_j b_je^{i\omega_jx}$ with
$|h\omega_j|<\pi$, the pole contribution on the existing midpoint contour
satisfies
\begin{equation}
 E_{\rm alias}:=\sup_{x\in\Omega}|\mathcal P_h[F_x]|
 \le\frac{2}{1-e^{-A_\lambda}}
       \sum_{\omega_j\ne0}|b_j|\mathcal T(\lambda,h|\omega_j|).
 \label{eq:bandwidth-weighted-bound}
\end{equation}
\end{proposition}
\begin{proof}
For a unit mode with $\omega>0$, the density $a_\gamma(z)=[f(z+id)-f(z-id)]/(2id)$ gives
$a_\gamma(z)=i\sinh(t_\omega)e^{i\omega z}/d$. At height
$(2\ell+1)d$, the density magnitudes at the four poles above and below
$x,x_0$ sum to $4\sinh(t_\omega)\cosh((2\ell+1)t_\omega)/d$.
Use the residues $1/(2\gamma)$ in absolute value and
the midpoint quadrature factor bound $(e^{(2\ell+1)A_\lambda}-1)^{-1}$
(Javed and Trefethen, 2014, Sections 4--5), to obtain
\begin{align*}
 E_{\rm alias}
 &\le8\sinh t_\omega\sum_{\ell\ge0}
       \frac{\cosh((2\ell+1)t_\omega)}{e^{(2\ell+1)A_\lambda}-1}\\
 &=2\sinh t_\omega\sum_{m\ge1}
       \left[\frac1{\sinh(mA_\lambda-t_\omega)}
             +\frac1{\sinh(mA_\lambda+t_\omega)}\right].
\end{align*}
Extending the positive residue sum and expanding its geometric denominators
gives the second line. Each successive $m$-term is at most
$e^{-A_\lambda}$ times its predecessor because
$\sinh(v+A_\lambda)\ge e^{A_\lambda}\sinh v$ for $v>0$.
Summation proves the unit-mode bound. Negative frequencies have the same
absolute bound, constants have zero density, and linearity gives
\qeqref{eq:bandwidth-weighted-bound}. The factor two accounts for anchoring
at $x_0$ as well as evaluation at $x$.
\end{proof}

\subsection{Compact score and scalar prescription}
Define the compact and small-angle scores
\begin{equation}
 \mathcal S(\lambda,N;\omega)=2e^{-A_\lambda}\sinh\chi_\omega,
 \qquad
 \mathcal S_{\rm sa}(\lambda,N;\omega)=2\chi_\omega e^{-A_\lambda}.
 \label{eq:bandwidth-scores}
\end{equation}
For $0<h\omega<\pi$,
\begin{equation}
 1\le\frac{\mathcal T}{\mathcal S}
 \le\frac1{1-e^{-(2A_\lambda-\chi_\omega)}},
 \qquad
 \frac{\mathcal S}{\mathcal S_{\rm sa}}
 =\frac{\sinh\chi_\omega}{\chi_\omega}.
 \label{eq:bandwidth-score-comparison}
\end{equation}
Indeed, expand $1/\sinh z=2e^{-z}/(1-e^{-2z})$ in the two terms of
\qeqref{eq:bandwidth-pair}. Their leading exponentials sum to
$4e^{-A_\lambda}\sinh t_\omega\cosh t_\omega=\mathcal S$, and the
smaller denominator is $1-e^{-(2A_\lambda-\chi_\omega)}$.
The compact approximation is therefore controlled throughout this band;
small-angle linearization additionally requires $\chi_\omega\ll1$.
The general-activation formula above specializes to this same first-pair
score; Figure~\ref{fig:bw-width-rules} plots its bandwidth roots.

For a representative frequency $\bar\omega$ and a dimensionless score budget
$\varepsilon_{\rm eff}>0$, the practical rule solves
\begin{equation}
 \boxed{2\exp\!\left(-\frac{\pi^2}{\lambda_{\rm pred}}\right)
 \sinh\!\left(\frac{2\pi\bar\omega}{N\lambda_{\rm pred}}\right)
 =\varepsilon_{\rm eff}.}
 \label{eq:bandwidth-prescription}
\end{equation}
Use an admissible bandwidth bracket on an increasing branch. In fact,
\begin{equation}
 \frac{\partial\log\mathcal S}{\partial\log\lambda}
 =A_\lambda-\chi_\omega\coth\chi_\omega>0
 \quad\hbox{if }A_\lambda-\chi_\omega\ge1,
 \label{eq:bandwidth-monotonicity}
\end{equation}
because $\chi\coth\chi<\chi+1$. Check the condition at the largest bracket endpoint and use bisection with
$\log\mathcal S=-(A_\lambda-\chi_\omega)
 +\log(-\operatorname{expm1}(-2\chi_\omega))$.
If the whole bracket satisfies the budget, select its largest endpoint;
if none does, report infeasibility. A zero frequency imposes no score constraint.

\subsection{Frequency input and tolerance}
For known coefficients use
$\bar\omega=\sum_j|b_j||\omega_j|/\sum_j|b_j|$.
For period-two sample data, take $m_F=M-1=sN$ distinct observations
$y_\ell=f(-1+2\ell/m_F)$, excluding the repeated endpoint from the transform,
and define
\begin{equation}
 \widehat b_k=\frac1{m_F}\sum_{\ell=0}^{m_F-1}
      y_\ell e^{-2\pi i k\ell/m_F},\qquad
 \widehat{\bar\omega}
 =\frac{\sum_{k\in\mathcal K}|\widehat b_k|\,|\pi k|}
        {\sum_{k\in\mathcal K}|\widehat b_k|},\qquad
 \mathcal K=\{-\lfloor m_F/2\rfloor,\ldots,\lceil m_F/2\rceil-1\}.
 \label{eq:bandwidth-frequency-estimate}
\end{equation}
Weights are amplitudes, not powers; the constant mode is included in the
denominator. A zero denominator returns no estimate. For nonperiodic targets,
the endpoint treatment and any spectral filtering must be specified: the raw
periodic transform can reflect endpoint mismatch rather than interior
variation. $m_F$ controls observation resolution; $N$ controls hidden-grid replicas.

For tanh, amplitude weighting matches the leading small-angle bound, which is
proportional to $\sum_j|b_j||\omega_j|$. At finite angles, however, convexity
of $\sinh$ gives
$\bigl(\sum_j|b_j|\bigr)\mathcal S(\bar\omega)
 \le\sum_j|b_j|\mathcal S(|\omega_j|)$,
with $\lambda,N$ held fixed. The centroid substitution is therefore a
heuristic, not an upper-bound-preserving reduction. To certify aliasing, use
\qeqref{eq:bandwidth-weighted-bound} with all relevant frequencies, or the
spectrum-free bound the representation theorem's pole estimate.

For one nonzero absolute frequency with coefficient mass
$B_{\rm spec}=\sum_j|b_j|$, an absolute alias budget $\varepsilon_a$ gives
the leading normalization $\varepsilon_{\rm eff}\simeq
\varepsilon_a/(2B_{\rm spec})$. A bound retains the corrections in
\qeqref{eq:bandwidth-weighted-bound} and
\qeqref{eq:bandwidth-score-comparison}; a relative budget first multiplies by
$\|f\|_{2,\Omega}$. Allocate the budget alongside the other allowances in
the representation and numerical recovery analysis, rather than necessarily setting it to unit
roundoff. The prescription is an aliasing criterion, not a characterization
of the total-error minimizer.

\subsection{Width dependence and induced slopes}
Along an interior root of \qeqref{eq:bandwidth-prescription}, with frequency
and tolerance fixed, implicit differentiation in continuous $N$ gives
\begin{equation}
 \frac{d\log\lambda_{\rm pred}}{d\log N}
 =\frac{\chi\coth\chi}{A_\lambda-\chi\coth\chi},
 \qquad
 \frac{\gamma_{\rm pred}}W
 =\frac{\lambda_{\rm pred}}2
   \frac{N}{N+2\lceil a_H\sqrt N\rceil+1},
 \label{eq:bandwidth-slope-scaling}
\end{equation}
where $\omega=\bar\omega$ and $\lambda=\lambda_{\rm pred}$ in the
definitions above. Thus the relative bandwidth varies
slowly when $A_\lambda\gg\chi\coth\chi$. Fixed $\lambda$ gives
$\gamma/W\to\lambda/2$; more generally, uniform positive lower and upper
bounds on the selected bandwidth give $\gamma_{\rm pred}=\Theta(W)$ on
that regime. Figure~\ref{fig:bw-width-rules} displays the selected bandwidth
itself against $N$. Its slowly varying large-width curves should
not be read as a fixed-precision convergence theorem for the refined rule.

\begin{figure}[tbp]
 \centering
 \includegraphics[width=\linewidth]{\BandwidthWidthFigure}
 \caption{Predicted bandwidth versus $N$ for tanh, exact
 GELU, and SiLU, with the same vertical scale in all three panels.
 Colored curves evaluate the frequency-dependent rule for
 the eight stated target functions at $\varepsilon_{\rm eff}=2^{-52}$.
 Each horizontal dashed line shows the activation-specific general
 Fourier-ratio root (approximately $0.25$ for tanh). These are scalar rule calculations, not
 network fits. The general rule is constant in width; the refined rule
 retains its frequency-dependent drift. Frequency conventions are stated
 above.}
 \label{fig:bw-width-rules}
\end{figure}

\begin{figure}[tbp]
 \centering
 \includegraphics[width=\linewidth]{\BandwidthErrorFigure}
 \caption{Measured relative error against bandwidth, compared with the
 general and frequency-dependent prescriptions. Columns: tanh, exact GELU,
 SiLU, and erf. Rows: the four targets written above each panel, all distinct
 from the eight targets of Figure~\ref{fig:bw-width-rules}. Colors show
 $N=32,48,64,96,128,256$ with 32 halo nodes per side. Dashed vertical lines
 mark the general-rule bandwidth; red diamonds mark the refined selections.
 Marker heights are directly evaluated network errors, not error
 predictions from the selector. All fits and evaluations use FP64;
 horizontal ranges differ between activations. The markers use the
 representative-frequency heuristic and need not minimize measured error.}
 \label{fig:bw-error-rules}
\end{figure}

\end{document}
